\section{Experimental Setup}
\label{sec:setup}

In this section we present our experimental setup for assessing the
performance of our \textbf{T}opic-aware \textbf{Conv}olutional
\textbf{S}equence \textbf{to} \textbf{S}equence model which we call
\textsc{T-ConvS2S} for short.  We discuss implementation details and
present the systems used for comparison with our approach.



\paragraph{Comparison Systems} 
We report results with various systems which were all trained on the
XSum dataset to generate a one-line summary given an input news
article.  We compared \mbox{\textsc{T-ConvS2S}} against three
extractive systems: a baseline which randomly selects a sentence from
the input document (\textsc{random}), a baseline which simply selects
the leading sentence from the document (\textsc{lead}), and an oracle
which selects a single-best sentence in each document
(\textsc{ext-oracle}). The latter is often used as an upper bound for
extractive methods. We also compared our model against the RNN-based
abstractive systems introduced by \newcite{see-acl17}. In particular,
we experimented with an attention-based sequence to sequence model
(\textsc{Seq2Seq}), a pointer-generator model which allows to copy
words from the source text (\textsc{PtGen}), and a pointer-generator
model with a coverage mechanism to keep track of words that have been
summarized (\textsc{PtGen-Covg}). Finally, we compared our model
against the vanilla convolution sequence to sequence model
(\textsc{ConvS2S}) of \newcite{convseq2seq}.



\paragraph{Model Parameters and Optimization} 

We did not anonymize entities but worked on a lowercased version of
the XSum dataset. During training and at test time the input document
was truncated to 400~tokens and  the length of the summary limited to
90~tokens.

The LDA model \cite{Blei:2003:LDA} was trained on XSum documents
(training portion). We therefore obtained for each word a probability
distribution over topics which we used to estimate $\mathbf{t'}$; the
topic distribution $t_D$ can be inferred for any new document, at
training and test time. We explored several LDA configurations on
held-out data, and obtained best results with~512
topics. Table~\ref{fig:lda-topics} shows some of the topics learned by
the LDA model.

%Once the LDA model is trained,
%we obtain a probability distribution over topics for each word in the
%training documents which we use to estimate $\mathbf{t'}$ and can
%infer the topic distribution $t_D$ for any new document, during both
%the training and test times.

For \textsc{Seq2Seq}, \textsc{PtGen} and \textsc{PtGen-Covg}, we used
the best settings reported on the CNN and DailyMail data
\cite{see-acl17}.\footnote{We used the code available at
  \url{https://github.com/abisee/pointer-generator}.}  All three
models had 256~dimensional hidden states and 128~dimensional word
embeddings. They were trained using Adagrad \cite{Duchi:2011} with
learning rate~0.15 and an initial accumulator value of~0.1. We used
gradient clipping with a maximum gradient norm of~2, but did not use
any form of regularization.  We used the loss on the validation set to
implement early stopping. 
%\textsc{Seq2Seq} takes around 4 days to
%train whereas \textsc{PtGen} and \textsc{PtGen-Covg} take around 2
%days each.\shaycomment{I am concerned that reviewers will claim that
%  you optimized hyperparameters on XSum for your model, but didn't do
%  the same for See's model.}

\begin{table}[t!]
  \center{\fontsize{10}{12}\selectfont % \footnotesize
    \begin{tabular}{@{}l@{\hspace{1.9ex}}p{6.7cm} }
      \thickhline 
      
      {T1:}& charge, court, murder, police, arrest, guilty, sentence, boy, bail, space, crown, trial \\
      
      {T2:}& church, abuse, bishop, child, catholic, gay, pope, school, christian, priest, cardinal \\
      
      {T3:}& council, people, government, local, housing, home, house, property, city, plan, authority \\
      
      {T4:}& clinton, party, trump, climate, poll, vote, plaid, election, debate, change, candidate, campaign\\
      
      {T5:}& country, growth, report, business, export, fall, bank, security, economy, rise, global, inflation\\
      
      {T6:}& hospital, patient, trust, nhs, people, care, health, service, staff, report, review, system, child\\

      \thickhline
    \end{tabular}     
  }
  \caption{Example topics learned by the LDA model on XSum documents
(training portion).}\label{fig:lda-topics}
\end{table}

For \textsc{ConvS2S}\footnote{We used the code available at
  \url{https://github.com/facebookresearch/fairseq-py}.} and
\textsc{T-ConvS2S}, we used 512 dimensional hidden states and 512
dimensional word and position embeddings. We trained our convolutional
models with Nesterov's accelerated gradient method
\cite{Sutskever:2013} using a momentum value of 0.99 and renormalized
gradients if their norm exceeded~0.1 \cite{Pascanu:2013}. We used a
learning rate of 0.10 and once the validation perplexity stopped
improving, we reduced the learning rate by an order of magnitude after
each epoch until it fell below~$10^{-4}$. We also applied a dropout
of~0.2 to the embeddings, the decoder outputs and the input of the
convolutional blocks. Gradients were normalized by the number of
non-padding tokens per mini-batch. We also used weight normalization
for all layers except for lookup tables.

All neural models, including ours and those based on RNNs \cite{see-acl17}
had a vocabulary of 50,000 words and were trained on a single Nvidia
M40 GPU with a batch size of~32 sentences.  Summaries at test time
were obtained using beam search (with beam size 10).

\section{Results}
\label{sec:results}


\paragraph{Automatic Evaluation} 

We report results using automatic metrics in Table~\ref{tab:rouge}.
We evaluated summarization quality using F$_1$ $\mbox{ROUGE}$
\cite{rouge}. Unigram and bigram overlap ($\mbox{ROUGE-1}$ and
$\mbox{ROUGE-2}$) are a proxy for assessing informativeness and the
longest common subsequence ($\mbox{ROUGE-L}$) represents
fluency.\footnote{We used \texttt{pyrouge} to compute all ROUGE
  scores, with parameters ``-a -c 95 -m -n 4 -w 1.2.''}



On the XSum dataset, \textsc{Seq2Seq} outperforms the \textsc{lead}
and \textsc{random} baselines by a large margin. \textsc{PtGen}, a
\textsc{Seq2Seq} model with a ``copying'' mechanism outperforms
\textsc{ext-oracle}, a ``perfect'' extractive system on ROUGE-2 and
ROUGE-L. This is in sharp contrast to the performance of these models
on CNN/DailyMail \cite{see-acl17} and  Newsroom datasets
\cite{newsroom-naacl18}, where they fail to outperform the
\textsc{lead}. The result provides further evidence that XSum is a
good testbed for abstractive summarization. \textsc{PtGen-Covg}, the
best performing abstractive system on the CNN/DailyMail datasets, does
not do well.  We believe that the coverage mechanism is more useful
when generating multi-line summaries and is basically redundant for
extreme summarization. 

\begin{table}[t]
  \begin{center}{\fontsize{8.5}{11}\selectfont %\footnotesize
  \begin{tabular}{ l | c c c }
    \thickhline
    Models & R1 & R2 & RL \\ \thickhline
    % \multicolumn{10}{c}{\textbf{Test Set (11334)}} \\ \hline
    Random & 15.16 & 1.78 & 11.27 \\ 
    \textsc{lead} & 16.30 & 1.60 & 11.95 \\ 
    \textsc{ext-oracle} & 29.79 & 8.81 & 22.66 \\ \hline
    \textsc{Seq2Seq} &  28.42 & 8.77 &  22.48 \\
    \textsc{PtGen} & 29.70 & 9.21 & 23.24 \\ 
    \textsc{PtGen-Covg} & 28.10 & 8.02 & 21.72 \\ 
    % \textsc{Transformer} & 28.63 & 9.58 & 23.29 \\ 
    \hline
    \textsc{ConvS2S} & 31.27 & 11.07 & 25.23 \\  % ppl drops to 23.4
    \textsc{T-ConvS2S} (enc$_{t'}$) & 31.71 & 11.38 & 25.56 \\
    \textsc{T-ConvS2S} (enc$_{t'}$, dec$_{t_D}$) & 31.71 & 11.34 & 25.61 \\
    \textsc{T-ConvS2S} (enc$_{(t',t_D)}$) & 31.61 & 11.30 & 25.51 \\
    \textsc{T-ConvS2S} (enc$_{(t',t_D)}$, dec$_{t_D}$)  & \textbf{31.89} & \textbf{11.54} & \textbf{25.75} \\ \thickhline % ppl drops to 23.04
    
    % \multicolumn{4}{c}{Weak Upper Bound} \\ \hline
    % Title  & 35.8 & 13.6 & 28.7 \\ \hline
  \end{tabular}
  }\end{center}
\caption{ROUGE results on  XSum test set. We report \mbox{ROUGE-1} (R1), ROUGE-2
  (R2), and ROUGE-L (RL) F$_1$  scores. Extractive systems are in the
  upper block, RNN-based   abstractive systems are in the middle block,
  and convolutional    abstractive systems are in the bottom
  block. \label{tab:rouge}} 
\end{table}

\begin{table}
  \begin{center}{\fontsize{8.5}{11}\selectfont % \footnotesize
  \begin{tabular}{ l | c c c c } 
    \thickhline
    \multirow{2}{*}{Models} & \multicolumn{4}{c}{\% of novel n-grams in generated summaries}  \\
    & unigrams & bigrams & trigrams & 4-grams  \\ \thickhline 
    \textsc{lead} & 0.00 & 0.00 & 0.00 & 0.00 \\
    \textsc{ext-oracle} & 0.00 & 0.00 & 0.00 & 0.00 \\
    % \textsc{Seq2Seq} & 36.66 & 82.17 & 95.58 & 98.63  \\
    \textsc{PtGen} & 27.40 & 73.33 & 90.43 & 96.04 \\ 
    % \textsc{PtGen-Covg} & 25.71 & 70.76 & 88.87 & 95.24 \\ \hline
    \textsc{ConvS2S} & \textbf{31.26} & \textbf{79.50} & \textbf{94.28} & \textbf{98.10}\\
    \textsc{T-ConvS2S}  & 30.73 & 79.18 & 94.10 & 98.03 \\ \hline
    \textsc{gold} & 35.76 & 83.45  & 95.50  & 98.49 \\ \thickhline
  \end{tabular}}
  \end{center}
  \caption{Proportion of novel $n$-grams in summaries generated by various models on the XSum test set.
    \label{tab:novelngram-xsum-models}}
\end{table}

\begin{table*}[t!]
  \center{ \footnotesize
    \begin{tabular}{l p{10.5cm} l}
      \thickhline 
      
      % 4306
      \textsc{ext-oracle} &Caroline Pidgeon is the Lib Dem candidate, Sian Berry will contest the election for the Greens and UKIP has chosen its culture spokesman Peter Whittle. &[34.1, 20.5, 34.1]\\
            
      {\textsc{PtGen}} & UKIP leader \textcolor{red}{Nigel Goldsmith} has been \textcolor{red}{elected} as the \textcolor{red}{new mayor of London} to elect a new conservative MP. &[45.7, 6.1, 28.6]  \\
      
      {\textsc{ConvS2S}}&  London mayoral candidate \textcolor{midnightblue}{Zac Goldsmith} has been \textcolor{red}{elected} as the \textcolor{red}{new mayor of London}.& [53.3, 21.4, 26.7]  \\

      {\textsc{T-ConvS2S}}& Former London mayoral candidate \textcolor{midnightblue}{Zac Goldsmith} has been \textcolor{midnightblue}{chosen to stand} in \textcolor{midnightblue}{the London mayoral election}. &[50.0, 26.7, 37.5] \\
      
      {\textsc{gold}}& \textcolor{midnightblue}{Zac Goldsmith} will \textcolor{midnightblue}{contest} the 2016 \textcolor{midnightblue}{London mayoral election} for the conservatives, it has been announced. \\
      
      {Questions}& (1) {Who will \textcolor{midnightblue}{contest} for
        the conservatives?} (\textcolor{midnightblue}{Zac
        Goldsmith})\\
 &  (2) {For what election will he/she contest?} (\textcolor{midnightblue}{The London mayoral election}) \\
            
      \hline
      % 4385
      {\textsc{ext-oracle}} &North-east rivals Newcastle are the only team below them in the Premier League table. &[35.3, 18.8, 35.3] \\
            
      {\textsc{PtGen}} &Sunderland have \textcolor{red}{appointed} former \textcolor{midnightblue}{Sunderland boss} \textcolor{midnightblue}{Dick Advocaat} as manager at the end of the season to sign a new deal. &[45.0, 10.5, 30.0] \\
      
      {\textsc{ConvS2S}} &Sunderland have \textcolor{red}{sacked} \textcolor{midnightblue}{manager} \textcolor{midnightblue}{Dick Advocaat} after less than three months in charge. &[25.0, 6.7, 18.8] \\

      {\textsc{T-ConvS2S}}& \textcolor{midnightblue}{Dick Advocaat} has \textcolor{midnightblue}{resigned} as \textcolor{midnightblue}{Sunderland manager} until the end of the season.& [56.3, 33.3, 56.3] \\
      
      {\textsc{gold}} &\textcolor{midnightblue}{Dick Advocaat} has \textcolor{midnightblue}{resigned} as \textcolor{midnightblue}{Sunderland boss}, with the team yet to win in the Premier League this season. \\
      
      Questions& (1) {Who has \textcolor{midnightblue}{resigned}?}
      (\textcolor{midnightblue}{Dick Advocaat}) \\& (2) {From what post has he/she resigned?} (\textcolor{midnightblue}{Sunderland boss}) \\
            
      \hline
      
      % 9590 
      {\textsc{ext-oracle}}& The Greater Ardoyne residents collective (GARC) is protesting against an agreement aimed at resolving a long-running dispute in the area. &[26.7, 9.3, 22.2]  \\

      {\textsc{PtGen}} &A residents' group has been granted permission for GARC to hold a parade on \textcolor{red}{the outskirts of Crumlin, County Antrim}. &[28.6, 5.0, 28.6] \\

      {\textsc{ConvS2S}} &A protest has been held in the \textcolor{midnightblue}{Republic of Ireland} calling for \textcolor{red}{an end to parading} parading in \textcolor{midnightblue}{North Belfast}. &[42.9, 20.0, 33.3] \\

      {\textsc{T-ConvS2S}}& A protest has been held in \textcolor{midnightblue}{North Belfast} over a protest against \textcolor{midnightblue}{the Orange Order} in \textcolor{midnightblue}{North Belfast}.& [45.0, 26.3, 45.0] \\

      {\textsc{gold}} &Church leaders have appealed to a nationalist residents' group to call off a protest against \textcolor{midnightblue}{an Orange Order parade} in \textcolor{midnightblue}{North Belfast}. \\

      Questions &(1) {Where is the protest supposed to happen?}
      (\textcolor{midnightblue}{North Belfast}) \\ 
&(2) {What are they protesting against?} (\textcolor{midnightblue}{An Orange Order parade}) \\
      \thickhline

    \end{tabular}     
  }
  \caption{Example output summaries on the XSum test set with [ROUGE-1,
    ROUGE-2 and ROUGE-L] scores, goldstandard reference, and
    corresponding questions. Words highlighted in blue are either the
    right answer or constitute appropriate context for inferring it;
    words in red lead to the wrong
    answer. }\label{tab:erroranalysis}
\end{table*}


\textsc{ConvS2S}, the convolutional variant of \textsc{Seq2Seq},
significantly outperforms all \mbox{RNN-based} abstractive systems. We
hypothesize that its superior performance stems from the ability to
better represent document content (i.e.,~by capturing long-range
dependencies).
% for abstractive summarization. \shaycomment{you keep
%  mentioning long-range dependencies, but do you ever show that this
%  is what we tackle?}
Table~\ref{tab:rouge} shows several variants of \textsc{T-ConvS2S}
including an encoder network enriched with information about how
topical a word is on its own (enc$_{t'}$) or in the document
(enc$_{(t',t_D)}$).  We also experimented with various decoders by
conditioning every prediction on the topic of the document, basically
encouraging the summary to be in the same theme as the document
(dec$_{t_D}$) or letting the decoder decide the theme of the
summary. Interestingly, all four \textsc{T-ConvS2S} variants
outperform \textsc{ConvS2S}. \textsc{T-ConvS2S} performs best when
both encoder and decoder are constrained by the document topic
(enc$_{(t',t_D)}$,dec$_{t_D}$). In the remainder of the paper, we
refer to this variant as \textsc{T-ConvS2S}.

We further assessed the extent to which various models are able to
perform rewriting by generating genuinely abstractive
summaries. Table~\ref{tab:novelngram-xsum-models} shows the proportion
of novel $n$-grams for \textsc{lead}, \textsc{ext-oracle},
\textsc{PtGen}, \textsc{ConvS2S}, and \textsc{T-ConvS2S}.  As can be
seen, the convolutional models exhibit the highest proportion of novel
$n$-grams. We should also point out that the summaries being evaluated
have on average comparable lengths; the summaries generated by
\textsc{PtGen} contain 22.57 words, those generated by
\textsc{ConvS2S} and \textsc{T-ConvS2S} have 20.07 and 20.22 words,
respectively, while \textsc{gold} summaries are the longest with~23.26
words.  Interestingly, \textsc{PtGen} trained on XSum only copies 4\%
of 4-grams in the source document, 10\% of trigrams, 27\% of bigrams,
and 73\% of unigrams. This is in sharp contrast to \textsc{PtGen}
trained on CNN/DailyMail exhibiting mostly extractive patterns; it
copies more than 85\% of 4-grams in the source document, 90\% of
trigrams, 95\% of bigrams, and 99\% of unigrams \cite{see-acl17}. This
result further strengthens our hypothesis that XSum is a good testbed
for abstractive methods.

% is able to calibrate its document
%modeling with the pertinent content in the document for abstractive
%summarization and the decoder (dec$_{t_D}$) is guided towards
%generating summaries on the same topic as the document. 


\paragraph{Human Evaluation}
In addition to automatic evaluation using ROUGE which can be
misleading when used as the only means to assess the informativeness
of summaries \cite{schluter:2017:EACLshort}, we also evaluated system
output by eliciting human judgments in two ways. 

In our first experiment, participants were asked to compare summaries
produced from the \textsc{ext-oracle} baseline, \textsc{PtGen}, the
best performing system of \citet{see-acl17}, \textsc{ConvS2S}, our
topic-aware model \textsc{T-ConvS2S}, and the human-authored gold
summary (\textsc{gold}). We did not include extracts from the
\textsc{lead} as they were significantly inferior to other models. 


The study was conducted on the Amazon Mechanical Turk platform using
\textit{Best-Worst Scaling} (BWS; \citealt{louviere1991best,
  louviere2015best}), a less labor-intensive alternative to paired
comparisons that has been shown to produce more reliable results than
rating scales \cite{bestworstscaling}. Participants were presented
with a document and summaries generated from two out of five systems
and were asked to decide which summary was \textit{better} and which one
was \textit{worse} in order of informativeness (does the summary
capture important information in the document?) and fluency (is the
summary written in well-formed English?). Examples of system summaries
are shown in Table~\ref{tab:erroranalysis}.  We randomly selected
50~documents from the XSum test set and compared all possible
combinations of two out of five systems for each document. We
collected judgments from three different participants for each
comparison. The order of summaries was randomized per document and the
order of documents per participant.

\begin{table}[t]
  \center{\small
    \begin{tabular}{lrc}
      \thickhline
      Models  & Score & QA\\ \hline
      \textsc{ext-oracle} &  -0.121 & 15.70\\
      \textsc{PtGen} & -0.218& 21.40 \\
      \textsc{ConvS2S}  &
      -0.130 & 30.90 \\
      \textsc{T-ConvS2S} & \textbf{0.037} & \textbf{46.05}\\ \hline
      \textsc{gold} & 0.431 & 97.23\\
      \thickhline
    \end{tabular}}
  \caption{System ranking according to human judgments  and QA-based
    evaluation.\label{tab:heval-pref}}
\end{table}

The score of a system was computed as the percentage of times it was
chosen as best minus the percentage of times it was selected as
worst. The scores range from -1 (worst) to 1 (best) and are shown in
Table~\ref{tab:heval-pref}. Perhaps unsurprisingly human-authored
summaries were considered best, whereas, \textsc{T-ConvS2S} was ranked
2nd followed by \textsc{ext-oracle} and
\textsc{ConvS2S}. \textsc{PtGen} was ranked worst with the lowest
score of $-0.218$. We carried out pairwise comparisons between all
models to assess whether system differences are statistically
significant. \textsc{gold} is significantly different from all other
systems and \textsc{T-ConvS2S} is significantly different from
\textsc{ConvS2S} and \textsc{PtGen} (using a one-way ANOVA with
posthoc Tukey HSD tests; $p < 0.01$). All other differences are not
statistically significant.


%\shaycomment{make sure all tense
%  is consistent in the experiments section. you say sometimes was
%  ranked, then is ranked, etc.}

%\begin{table}[t]
%  \center{\footnotesize
%    \begin{tabular}{|@{~}l@{~~}| c@{~~} c@{~~} c@{~~} c@{~~} |c@{~}| } 
%      \hline
%      Models & \textsc{\scriptsize ext-oracle} & \textsc{\scriptsize PtGen} & \textsc{\scriptsize  ConvS2S} & \textsc{\scriptsize T-ConvS2S} & \textsc{\scriptsize gold} \\ \hline
%      QA & 15.70 & 21.40 & 30.90 & \textbf{46.05} & 97.23 \\
%      \hline
%    \end{tabular}}
%  \caption{QA-based evaluation. The row QA shows the percentage of questions
%    that participants answered correctly by reading system summaries.\label{tab:heval-qa}}
%\end{table}
%



For our second experiment we used a question-answering (QA) paradigm
\cite{Clarke:Lapata:2010,narayan-rank18} to assess the degree to which
the models retain key information from the document. We used the same
50 documents as in our first elicitation study.  We wrote two
fact-based questions per document, just by reading the summary, under
the assumption that it highlights the most important content of the
news article. Questions were formulated so as not to reveal answers to
subsequent questions.  We created 100~questions in total (see
Table~\ref{tab:erroranalysis} for examples). Participants read the
output summaries and answered the questions as best they could without
access to the document or the gold summary. The more questions can be
answered, the better the corresponding system is at summarizing the
document as a whole.  Five participants answered questions for each
summary.

We followed the scoring mechanism introduced in
\newcite{Clarke:Lapata:2010}. A correct answer was marked with a score
of one, partially correct answers with a score of~0.5, and zero
otherwise. The final score for a system is the average of all its
question scores. Answers again were elicited using Amazon's Mechanical
Turk crowdsourcing platform. We uploaded the data in batches (one
system at a time) to ensure that the same participant does not
evaluate summaries from different systems on the same set of
questions.

%Example questions are given in
%Figure~\ref{fig:erroranalysis}.

% Example questions are given in Figure~\ref{fig:erroranalysis}. Subjects
% were shown summaries from three systems: the \textsc{Lead} baseline,
% the abstractive system of \newcite{see-acl17}, and \refresh.




%\paragraph{Human Evaluation: System Ranking}

 
%\paragraph{Human Evaluation: Question Answering}

Table~\ref{tab:heval-pref} shows the results of the QA
evaluation. Based on summaries generated by \mbox{\textsc{T-ConvS2S}},
participants can answer $46.05\%$ of the questions
correctly. Summaries generated by \textsc{ConvS2S}, \textsc{PtGen} and
\textsc{ext-oracle} provide answers to $30.90\%$, $21.40\%$, and
$15.70\%$ of the questions, respectively. Pairwise differences between
systems are all statistically significant ($p < 0.01$) with the
exception of \textsc{PtGen} and \textsc{ext-oracle}.
\textsc{ext-oracle} performs poorly on both QA and rating
evaluations. The examples in Table~\ref{tab:erroranalysis} indicate
that \textsc{ext-oracle} is often misled by selecting a sentence with
the highest ROUGE (against the gold summary), but ROUGE itself does
not ensure that the summary retains the most important information
from the document.  The QA evaluation further emphasizes that in order
for the summary to be felicitous, information needs to be embedded in
the appropriate context. For example, \textsc{ConvS2S} and
\textsc{PtGen} will fail to answer the question ``Who has resigned?''
(see Table~\ref{tab:erroranalysis} second block) despite containing
the correct answer ``Dick Advocaat'' due to the wrong
context. \mbox{\textsc{T-ConvS2S}} is able to extract
important entities from the document with the right theme.

%Our analysis suggests that ROUGE should
%not be the only criteria to evaluate summarization systems, it should
%always be complemented with human evaluations.

%The question answering evaluation not only judges if the summary
%retains the main information from the document, but it also 



% A: "LEAD",
% B: "EXT-Oracle"
% C: "PTGen"
% D: "ConvS2S"
% E: "T-ConvS2S"
% 6: "GOLD"

%         Tukey HSD
% Q statistic        Tukey HSD
% p-value        Tukey HSD
% inferfence
% A vs B        1.7924        0.7769044        insignificant
% A vs C        5.1981        0.0032943        ** p<0.01
% A vs D        10.8742        0.0010053        ** p<0.01
% A vs E        19.5974        0.0010053        ** p<0.01
% A vs F        50.5470        0.0010053        ** p<0.01
% B vs C        3.4056        0.1536003        insignificant
% B vs D        9.0817        0.0010053        ** p<0.01
% B vs E        17.8050        0.0010053        ** p<0.01
% B vs F        48.7545        0.0010053        ** p<0.01
% C vs D        5.6761        0.0010053        ** p<0.01
% C vs E        14.3993        0.0010053        ** p<0.01
% C vs F        45.3489        0.0010053        ** p<0.01
% D vs E        8.7232        0.0010053        ** p<0.01
% D vs F        39.6728        0.0010053        ** p<0.01
% E vs F        30.9496        0.0010053        ** p<0.01




%\shashi{Anything else?}

%% We also found that ConS2s and TCon sec seems to get entities right. 
%% rEPEATITION, PTGEN IN CONV COULD HELP, WITHPOUT Copy MECHANISM


% May be T-Conv more factually correct than other. See this. Find other example.

% Abs methods might have correct answer but some wrong info, that is then causes difficulties to answer by users. (Find an example in support of this). However, ext always fluent, is found answer it will be correct.

% T-ConvS2s successfull achieves full name but fails to answer becuase misses the right context. But that is not the case with ext system.

% TConVseq or convs2s fails because does not get proper context: Despite of capturing the named entity. PTGen often find wron combo.

% [4306]Who will be the conservative contestant?
% 5 [1, 0, 0, 0, 0] T-ConvS2S 0.2
% 5 [0, 0, 0, 0, 0] LEAD 0.0
% 5 [1, 0, 1, 1, 1] GOLD 0.8
% 5 [0, 0, 0, 0, 0] PTGen 0.0
% 5 [0, 0, 0, 0, 0] EXT-Oracle 0.0
% 5 [0, 0, 0, 0, 0] ConvS2S 0.0

% [4385]Who has resigned?
% 5 [1, 1, 1, 1, 1] T-ConvS2S 1.0
% 5 [0, 0, 0, 0, 0] LEAD 0.0
% 5 [1, 1, 1, 1, 1] GOLD 1.0
% 5 [0, 0, 0, 0, 0] PTGen 0.0
% 5 [0, 0, 0, 0, 0] EXT-Oracle 0.0
% 5 [0, 1, 1, 0, 0] ConvS2S 0.4


%% Age etc is still a problem for all the model.

%% T-ConvS2S
%% the ska band , bradbury , has died at the age of 77 , the band has announced .
%% [1353]At what age did he die? A1GV0UZU0T2ORS
%% Candidate Answer: 77
%% Reference Answer: 62


%% Match reports
%% 2783   38782795 

% Position in RNN shown to learn by itlesf. \cite{D16-1248}
% Compare with seq2seq abssum models: We don't add position embed to seq2seq. No point, people have shown they already learn that in their neurons.


%%% Local Variables:
%%% mode: latex
%%% TeX-master: "extreme-sum"
%%% End:

% \begin{table}
%   \begin{center}{\footnotesize
%   \begin{tabular}{ l | c c c } 
%     % \multicolumn{4}{c}{\textsc{NewsRoom-Abstractive} Dataset} \\
%     \thickhline
%     Models & R1 & R2 & RL \\ \thickhline 
%     \multicolumn{4}{c}{\textsc{development set}} \\ \hline 
%     \textsc{lead-3} & 13.69 & 2.42 & 11.24 \\
%     % Fragments &  73.43 & 47.66 & 73.43 \\ 
%     TextRank & 13.54 & 1.88 & 10.46 \\
%     Abs-N & 6.21 & 1.07 & 5.68 \\
%     \textsc{PtGen} &  14.66 & 2.26 & 11.44 \\ \hline
%     \textsc{lead-2} & 15.41 & 2.73 & 12.28\\
%     \textsc{ext-oracle} & \textbf{29.77} & \textbf{7.83} & \textbf{24.81}\\
%     \textsc{ConvS2S} & 16.61 & 5.37 & 14.38 \\
%     \textsc{T-ConvS2S} & 16.88 & 5.44 & 14.57 \\ \thickhline 
%     \multicolumn{4}{c}{\textsc{test set}} \\ \hline 
%     % \textsc{lead-1} & 13.98 & 2.48 & 10.85\\
%     \textsc{lead-2} & 15.44 & 2.72 & 12.32 \\
%     \textsc{lead-3} & 14.99 & 2.69 & 12.17\\
%     \textsc{ext-oracle} & \textbf{29.85} & \textbf{7.82} & \textbf{24.86}\\
%     \textsc{ConvS2S} & 16.77 & 5.57 & 14.54\\
%     \textsc{T-ConvS2S} & 16.97 & 5.56 & 14.70\\ \thickhline 

%   \end{tabular}}
%   \end{center}
%   \caption{ROUGE results on Newsroom-Abs subset (development and released test sets). 
%     \protect\newcite{newsroom-naacl18} have reported the performances of TextRank \protect\cite{Mihalcea04TextRank}, Abs-N \protect\cite{rush-acl15} and \textsc{PtGen} \protect\cite{see-acl17} on the development set (first block) only, these results are not available on the test set. They also report the \textsc{lead-3} baseline of selecting first three sentences, we found that \textsc{lead-2} (selecting first two sentences) performs better than \textsc{lead-3}.}\label{tab:newsroom-results} 
% \end{table}

% \paragraph{Results on Newsroom Dataset}

% Table~\ref{tab:newsroom-results} presents the performance comparison
% on the abstractive subset of Newsroom. \textsc{ConvS2S} and
% \textsc{T-ConvS2S} significantly outperform the \textsc{lead}
% baselines on both the development and test sets. They also
% significantly outperform \textsc{PtGen} on the development set. The
% \textsc{PtGen} results on the test set are not made available by
% \newcite{newsroom-naacl18}. \textsc{T-ConvS2S} is superior to
% \textsc{ConvS2S} on both the development and test sets. This is in
% line with our observations on the XSum dataset. However, in contrast
% to the XSum dataset where the abstractive models \textsc{PtGen},
% \textsc{ConvS2S} and \textsc{T-ConvS2S} are able to outperform the
% extractive oracle (\textsc{ext-oracle}), on the Newsroom dataset the
% abstractive models are far behind \textsc{ext-oracle}. It suggests
% that improved sentence selection would bring further performance gains
% for extractive approaches on the Newsroom dataset.
